
\begin{titlepage}
% \begin{CJK}{UTF8}{titlepage}
%  \CJKfamily{mj}
\pagenumbering{gobble}
\begin{center}
% \vspace*{0.5cm}
\Huge
검증 가능한 생체분자 계산 워크플로우를 위한 블록체인 기반 신뢰 인프라


\vspace{1.5cm}
\LARGE

 Victor Ikenna Kanu
\vspace{1.5cm}
\LARGE

% Department of IT Convergence Engineering,\\
% Graduate School\\
% Kumoh National Institute of Technology
국립금오공과대학교 대학원 IT융복합공학과
\vspace{2cm}

% Abstract

요    약
\end{center}
\large

본 학위논문은 변조 탐지 가능한 출처 정보(provenance), 결정론적 출력 검증, 그리고 다자간 교환을 위한 정책 기반 유효성을 통합하는 검증 가능한 시스템으로서 생체분자 계산 워크플로우(BCWs)를 공학적으로 구현하는 방법을 연구한다. 본 연구에서는 네 가지 상호 보완적인 구현체를 설계하고, 공유 검증 패턴 하에서 허가형 원장 인프라(RQ1·RQ3·RQ4는 PureChain, RQ2는 Hyperledger Fabric)에 배포하였으며, 네 가지 연구 질문을 다루는 대표적인 생의학 파이프라인에서 실증적으로 평가하였다. RQ1(증거 및 출처 정보)에서는 허가형 원장, 분산 저장소, 역할 기반 접근 제어를 통합한 블록체인 기반 마이크로서비스 프레임워크가 485.96 req/s의 API 처리량, 445 TPS의 온체인 성능, 그리고 SHA-256 기반 변조 탐지에서 0건의 위양성을 달성하였으며, HIV-1 프로테아제 사례 연구를 통해 종단 간 출처 정보가 검증되었다. RQ2(검증 가능한 연산)에서는 3계층 블록체인 통합 검증 시스템이 출처 정보를 변조 탐지 이력에서 다이제스트 기반 검증으로 확장하여, 5개 단백질 표적에 대해 100\%의 재현성을, 6개 구조적 범주의 20개 표적에 대해 98.2\%의 재현성을 달성하였다. RQ3(실제 환경에서의 결정성)에서는 10개 표적에 걸친 71{,}853개 화합물을 처리하는 결정론적 블록체인 감사 가상 스크리닝 파이프라인이 우수한 예측 성능(평균 $\mathtt{R^2}$ = 0.693, AUC-ROC = 0.937)을 보이는 동시에, 40회의 독립 재실행에 걸쳐 100\% 해시 검증을 통과함으로써 모든 구성요소 수준의 결정성 검사를 통과하였다. RQ4(정책 기반 유효성)에서는 정책 기반 블록체인-자격증명 구조가 품질 관리 및 상태 이상 탐지에서 F1 = 1.00을 달성하였으며, 레지스트리 크기와 무관하게 p50 $\sim$ 32 ms의 검증 지연시간을 보였다. 종합적으로, 이 네 가지 상호 보완적인 구현체는 3차원 유효성 결합 Valid(w) = AuditValid(w) $\wedge$ DeterminismValid(w) $\wedge$ PolicyValid(w)를 실증적으로 뒷받침하며, 재현 가능하고 검증 가능하며 거버넌스 준수가 가능한 생체분자 계산 워크플로우를 위한 조합 가능한 신뢰 스택을 제공한다. 본 신뢰 스택에서 각 유효성 차원은 독립적으로 입증되며, 그 결합은 단일 통합 실행이 아닌 구조적 차원에서 특성화된다.


\vspace{0.5cm}
\textbf{키워드:} 생체분자 계산 워크플로우, 블록체인, 프로비넌스, 결정론, 검증 가능한 계산, 정책 집행.


% \end{CJK}
\end{titlepage}